% TOP_PROBLEM: {"id": 3700039, "title": "Median inequality for three subharmonic functions", "category_id": 8, "category": {"id": 8, "name": "analysis", "display_name": "Analysis", "description": "Limits, continuity, calculus, and function theory.", "slug": "analysis", "order_index": 8, "created_at": "2026-07-31T15:26:25.671Z"}, "status": "open", "problem_number": "AMR-036-0039", "set_id": 13, "statusQualification": "Corrects the source's rearrangement typo and defines the harmless 0/0 case explicitly."}
% FIRST_POSTED: 2026-10-01
% ENTRY_KIND: statement_only
% UnsolvedMath ID 3700039; source catalogue code AMR-036-0039
% !TeX program = lualatex
% Definition and sources only; this file is not an LLM solution attempt.
\documentclass[11pt,a4paper]{article}
\usepackage[margin=25mm]{geometry}
\usepackage{fontspec}
\setmainfont{lmroman10-regular.otf}[BoldFont=lmroman10-bold.otf,
 ItalicFont=lmroman10-italic.otf,BoldItalicFont=lmroman10-bolditalic.otf]
\usepackage{amsmath,amssymb,mathtools}
\usepackage{unicode-math}
\setmathfont{latinmodern-math.otf}
\AtBeginDocument{\let\setminus\smallsetminus}
\usepackage{xurl,hyperref}
\hypersetup{colorlinks=true,urlcolor=blue}
\setcounter{secnumdepth}{0}
\setlength{\parindent}{0pt}
\setlength{\parskip}{5pt}
\setlength{\emergencystretch}{3em}
\begin{document}
\section*{Median inequality for three subharmonic functions}\label{3700039}
{\small UnsolvedMath ID 3700039 · AMR-036-0039 · Analysis · AMR Open Problem Lists}
\subsection{Definitions and mathematical statement}
Let $u_1,u_2,u_3$ be subharmonic functions on
$\mathbb C$ with $u_j(0)=0$. Subharmonic means
upper semicontinuity together with the submean inequality
on every disk; values $-\infty$ away from the origin
are allowed. At each point arrange the three values
in increasing order, giving
$v_1(z)\leq v_2(z)\leq v_3(z)$.
Thus $v_2$ is their median and $v_3$ their maximum;
the median need not itself be subharmonic.

For $r>0$ put
\[
 I(r)=\int_0^{2\pi}v_2(re^{i\theta})\,d\theta,\qquad
 B(r)=\max_{|z|=r}v_3(z).
\]
The angular integral is unnormalized. The submean
inequality at the origin implies $B(r)\geq0$.
Define $R(r)=I(r)/B(r)$ when $B(r)>0$,
and set $R(r)=0$ when $B(r)=0$.
In the latter case the maximum principle forces
all three functions to vanish on the disk,
so $I(r)=0$ as well.

Must every such triple satisfy
$\sup_{r>0}R(r)\geq0$?
Does the stronger local assertion
$\limsup_{r\downarrow0}R(r)\geq0$ always hold?
These statements allow the ratio to remain
negative at individual radii, or even approach
zero from below; they do not demand a positive
median mean at some specified radius.
No growth restriction, harmonicity assumption,
or common homogeneity is imposed on the three
subharmonic functions.
\subsection{Short English statement}
Must the normalized circular mean of the median of three subharmonic functions be nonnegative in supremum, even arbitrarily near the origin?
\subsection{Sources}
\begin{itemize}
\item[S1] ULAM.AI and the UnsolvedMath Contributors, supplied problems.json, record 3700039 (AMR-036-0039), AMR Open Problem Lists. Catalogue identity and source transcription; CC BY 4.0.
\url{https://huggingface.co/datasets/ulamai/UnsolvedMath}.
\item[S2] Eremenko - My favorite unsolved problems, Subharmonic inequality conjecture. Original source indexed by AMR-036-0039.
\url{https://www.math.purdue.edu/~eremenko/uns1.html}.
\item[S3] A. Eremenko, Conjectured inequality for subharmonic functions. Author-maintained primary problem note.
\url{https://www.math.purdue.edu/~eremenko/dvi/subh1.pdf}.
\end{itemize}
\end{document}
